\honest{Many Worlds, One Best Guess}{Eliezer Yudkowsky}{2008}

Photons, electrons, atoms and lasers seem to be governed by simple laws. Which generalization of them should you believe? Call something grue if it looks green before 2020 and blue after. Just as ``Emeralds are green'' beats ``Emeralds are grue,'' the simplest one wins: that all electrons have spin 1/2, not that they do before 2020, or unless they are part of an entangled system heavier than a gram. The laws hold everywhere. \nb{``Grue'' is Nelson Goodman's puzzle. The whole post rests on this rule, and defers its defence.}

We cannot check the equations in detail for a human, but quantum effects do not vanish in bulk: lasers emit floods of coherent photons, atoms have the chemistry the theory says, molecules hold together, and every prediction tested has come true. Someone could claim a difference that shows up only ``at the unmeasurable 20th decimal place'' of small things, and ``we can't prove they're wrong.'' Occam's razor asks why we should think about that one new law among zillions.

Apply the laws to Schrödinger's cat, and the cat and the person who looks end up in superposition. The brain that sees the exploded cat and the one that sees the living cat differ in many firing neurons, so they are far apart in configuration space and interact only in about the $10^{30}$th decimal place; no mind sees a blurred cat. Seeing one outcome is exactly what the simplest laws predict, and superposition is being verified at ever larger scales, with a 50-micron test under way. \nb{This is the many-worlds answer to the ``preferred basis'' objection, which asks what counts as a world. The encyclopedia entry that argues for many-worlds says it is ``no longer considered a serious objection''; Jess Riedel, quoted later in the post, saw ``serious conceptual problems with defining individual worlds, even emergently.''}

So the other Earths are not an extra assumption. They are ``a logical consequence'' of the simplest generalization, and to call them unnecessary entities is ``just not how Occam's Razor works.'' \nb{They follow only if the wavefunction is all there is. Bohm's theory keeps the same wave equation everywhere, adds particle positions, has one world and predicts the same results. Counting laws rather than objects, the encyclopedia entry on many-worlds agrees that it is the most economical theory; critics reply that its laws alone do not explain experience.}

One puzzle remains: the Born probabilities. Counting observers seems to predict 50/50 for every two-way experiment, and that is not what we see. Robin Hanson suggests that worlds of exponentially small amplitude are disrupted by leakage from larger ones, which would give back the Born probabilities; I like it because it is so normal. My own thought: counting observers the obvious way gives strange results even outside quantum physics. If I split my brain into a trillion copies whenever I win a lottery while anaesthetized, then merge them, naive counting says I can make myself win. The Born rule has no such split-and-merge problem; given unitary physics it is the unique rule that stops observers having psychic powers, which is interesting though it explains nothing. I am suspicious of Wallace's and Deutsch's derivations from decision theory, because there seems to be a part of ``What happens to me?'' that my utility function cannot change: even if I cared nothing for the worlds where I lose a quantum lottery, I would still ``mostly'' wake up in one. \nb{Much of the published literature shares that suspicion. The usual many-worlds answer adds a postulate about probability, and how there can be probability when every outcome happens is called ``arguably the most serious difficulty'' for the view.} There is hope of an answer without new fundamental laws, since we split continuously, not in two at each reading of a detector, and the anthropic weight of observers or the brain's superpositions may hold surprises. But I cannot rule a new law out. Jess Riedel says every new experimental regime has brought new physics; ``Every time'' is too strong, but a law hidden in the 20th decimal place is possible.

Might a new law leave only one world? Then all our other selves are gone, and if the universe is also small, this could be the only Earth anywhere. But you should not even ask. Asking is the error of the Intelligent Design folk, who pick one explanation out of a zillion, and it ``betrays a sentimental attachment to human intuitions already proven wrong.'' Asking whether a new law might give the universe a centre at Earth, rather than at Proxima Centauri, or a favourite pizza topping of pepperoni, betrays the agenda; and a centre that follows Earth through space would be hard to give any simple law. ``The wheel of science turns, but it doesn't turn backward.'' \nb{The single-world theories are specific, worked-out rivals, not random picks, and their defenders give arguments, for instance that a wavefunction on configuration space cannot by itself describe objects in space; the post gives no evidence for the motive it names.}

I have two specific reasons against one world. First, I ask how a fundamental law could pick out one high-level world. \nb{The actual theories do not try: GRW localizes particles at random times, and Bohm gives each particle a position.} Second, and much worse, a single world ``would violate Special Relativity.'' With many worlds, when you measure one of an entangled pair of photons, each of you sees 50/50, and so does a friend a light-year away measuring at 20 degrees to your basis; only when you meet years later is the chance that you got the same result 11.6 per cent. With one world, the friend's probability actually changes from 50/50 to 11.6 per cent when you measure. By Bell's theorem the results cannot have been there all along, so something changes faster than light, in any single-world theory. I dismiss the backward-in-time escape. The experiments are done: any influence would have to travel six million times as fast as light, and it ignores which measurement comes first. \nb{Those figures check. The step from nonlocal to violating relativity does not hold for every theory. Bohm-type theories need a preferred slicing of spacetime, but for random collapse theories the correlation ``does not require that one of the events be in the past of the other,'' and fully relativistic versions of GRW were published by 2006.} Relativity is not an arbitrary speed limit that a crime could evade by going back in time; it says there is no global now, no before or after across spacelike gaps, and the information for a single global world is simply not present locally. My logic may contain ``some hidden assumption,'' I grant. Still, the single world was an ancient intuition, ``disproved'' like absolute time. \nb{The disproof is the argument just described, which the post has just allowed may be flawed; the larger superpositions it cites are predicted by Bohm's single-world theory too.} Quantum non-realism I call ``a Get Out of Jail Free card'': ``It's okay to violate Special Relativity because none of this is real anyway!'' \nb{That slogan is the post's. QBists, for example, deny that there is any nonlocal influence to begin with.}

It is ``not quite safe'' to call other Earths as well established as any other truth of science; they are less certain than trees. But one Earth is about as likely as a spinning black hole breaking conservation of angular momentum. We have deep reasons, from the rotational symmetry of space, to expect conservation everywhere, though as far as I know no one has checked black holes (or my flushing toilet); and if you dwell on that one possibility, especially without knowing why angular momentum is conserved, it starts to seem plausible. Its rational probability is small, and so is that of one Earth. Many-worlds is ``an obvious fact, if you have all your marbles lined up correctly.'' It is not universally accepted only because of a historical accident and because ``many academic physicists do not have a mathematical grasp of Occam's Razor.'' I am not in academia, so I speak ``on behalf of any and all physicists out there who dare not say it themselves'': many-worlds ``wins outright''; a single Earth is no more likely than colliding top quarks breaking conservation of energy, and needs not just an unknown law but magic. The debate should have been over fifty years ago; ``There is no rational controversy to teach''; our children will look back and correctly deduce that we were nuts. \nb{No evidence is given for these explanations of dissent. Specialists who know the formal arguments have published objections, and open defenders of many-worlds wrote for decades before the post. In polls from 1998 to 2025, 15 to 18 per cent of respondents chose many-worlds as their favourite.} Last: for the honour of my Earth, I write as if many-worlds ``were an established fact,'' because it is one; the only question is how long this world will take to update.
